% SEGMENT OLP-0467-S01
% Part: normal-modal-logic
% Chapter: tableaux
% Section: simple-S5

\documentclass[../../../include/open-logic-section]{subfiles}

\begin{document}

\olfileid{nml}{tab}{s5}

\olsection{\Log{S5} க்கான எளிய \usetoken{P}{tableau}கள்}

அனைத்துறவு மாதிரிகளின் வகுப்பை ஒப்பிட்டு~\Log{S5} சரியானதும் நிறைவானதும்;
அதாவது, இம்மாதிரிகளில் ஒவ்வோர் உலகும் ஒவ்வோர் உலகிலிருந்தும் அணுகத்தக்கது.
அனைத்துறவு மாதிரிகளில் அணுகுறவு பொருட்டல்ல: ``$\mSat{M}{!A}[w]$ ஆகும்
ஓர் உலகம்~$w$ உள்ளது'' என்பது,~$u$ இலிருந்து அணுகத்தக்க அத்தகைய ஒரு~$w$
உள்ளது என்பதற்கு நிகரானது. எனவே \Log{S5} இல், மாதிரிகளை உலகங்களின் ஒரு
கணமும் ஒரு மதிப்பீடு~$V$ உம் என்று எளிதாக வரையறுக்கலாம். !!{tableau}
விதிகளையும் நாம் எளிமைப்படுத்த இயலும் என்பதை இது உணர்த்துகிறது. பொது
நேர்வில், எந்த முன்னொட்டுகள் எந்த மற்ற முன்னொட்டுகள் பெயரிடும்
உலகங்களிலிருந்து அணுகத்தக்க உலகங்களுக்குப் பெயரிடுகின்றன என்பதைக் கண்காணிக்க,
நேர்ம முழுக்களின் தொடர்களை முன்னொட்டுகளாகக் கொள்கிறோம்:~$\sigma$ விலிருந்து
அணுகத்தக்க ஓர் உலகிற்கு $\sigma.n$ பெயரிடுகிறது. ஆனால் \Log{S5} இல் எந்த
உலகும் எந்த உலகிலிருந்தும் அணுகத்தக்கது; ஆகவே இதைக் கண்காணிக்க வேண்டியதில்லை.
அதற்குப் பதிலாக நேர்ம முழுக்களையே முன்னொட்டுகளாகப் பயன்படுத்தலாம்.
எளிமைப்படுத்தப்பட்ட விதிகள் \olref{tab:rules-S5} இல் கொடுக்கப்பட்டுள்ளன.

\begin{table}
  \begin{center}
    \def\arraystretch{3}\def\fCenter{}
    \begin{tabular}{|c|c|}
    \hline
    \iftag{prvBox}{
      \AxiomC{\sFmla{\True}{\Box !A}[n]}
      \RightLabel{\TRule{\True}{\Box}}
      \UnaryInfC{\sFmla{\True}{!A}[m]}
      \DisplayProof
      &
      \AxiomC{\sFmla{\False}{\Box !A}[n]}
      \RightLabel{\TRule{\False}{\Box}}
      \UnaryInfC{\sFmla{\False}{!A}[m]}
      \DisplayProof\\
      $m$ பயன்படுத்தப்பட்டது & $m$ புதியது
      \\[1ex]
      \hline}{}
    \iftag{prvDiamond}{
      \AxiomC{\sFmla{\True}{\Diamond !A}[n]}
      \RightLabel{\TRule{\True}{\Diamond}}
      \UnaryInfC{\sFmla{\True}{!A}[m]}
      \DisplayProof
      &
      \AxiomC{\sFmla{\False}{\Diamond !A}[n]}
      \RightLabel{\TRule{\False}{\Diamond}}
      \UnaryInfC{\sFmla{\False}{!A}[m]}
      \DisplayProof\\
      $m$ புதியது & $m$ பயன்படுத்தப்பட்டது
      \\[1ex]
      \hline}{}
    \end{tabular}
  \end{center}
  \caption{\Log{S5} க்கான எளிமைப்படுத்தப்பட்ட விதிகள்.}
  \ollabel{tab:rules-S5}
\end{table}

\begin{ex}
  $\Log{S5} \Proves \Ax{5}$, அதாவது $\Diamond!A \lif
  \Box\Diamond!A$, என்பதைக் காட்டும் எளிமைப்படுத்தப்பட்ட மூடிய அட்டவணை
  மரத்தைக் கொடுக்கிறோம்.
  \begin{oltableau}
    [\pFmla{\False}{\Diamond\formula{A} \lif \Box\Diamond \formula{A}}{1},
      just = \TAss
      [\pFmla{\True}{\Diamond \formula{A}}{1}, just = {\TRule{\False}{\lif}[1]}
        [\pFmla{\False}{\Box\Diamond \formula{A}}{1},
          just = {\TRule{\False}{\lif}[1]}
          [\pFmla{\False}{\Diamond \formula{A}}{2},
            just = {\TRule{\False}{\Box}[3]}
            [\pFmla{\True}{\formula{A}}{3},
              just = {\TRule{\True}{\Diamond}[2]}
                [\pFmla{\False}{\formula{A}}{3},
                  just = {\TRule{\False}{\Diamond}[4]}, close]
            ]
          ]
        ]
      ]
    ]
  \end{oltableau}
\end{ex}

\end{document}
